\documentclass[11pt]{article}
% ===========================================================================
%  Shared preamble for the AP Calculus AB exam-preparation worksheets.
%  Input right after \documentclass. Files beginning with "_" are not compiled
%  on their own. Navy and gold, to match the course's lesson pages.
% ===========================================================================
\usepackage[letterpaper,top=2.3cm,bottom=2.2cm,left=1.7cm,right=1.7cm,headheight=28pt,footskip=30pt]{geometry}
\usepackage{amsmath,amssymb}
\usepackage[T1]{fontenc}
\usepackage[utf8]{inputenc}
\usepackage{lmodern}
\usepackage{xcolor}
\usepackage[most]{tcolorbox}
\usepackage{enumitem}
\usepackage{fancyhdr}
\usepackage{array}
\usepackage{tikz}
\usepackage{pgfplots}
\pgfplotsset{compat=1.18}

\definecolor{navy}{HTML}{161E54}
\definecolor{ink}{HTML}{27357E}
\definecolor{gold}{HTML}{B98A00}
\definecolor{goldbg}{HTML}{FFFDF3}
\definecolor{inkbg}{HTML}{F6F7FD}
\definecolor{rule}{HTML}{C9CFF0}
\definecolor{okgreen}{HTML}{1B7A44}
\definecolor{okbg}{HTML}{F1F8F3}

\pagestyle{fancy}
\fancyhf{}
\fancyhead[L]{\footnotesize Name: \underline{\hspace{4cm}}}
\fancyhead[C]{\textbf{\color{navy}Integration Academy}\\[-2pt]{\scriptsize Calculus AB \textperiodcentered{} AP\textsuperscript{\textregistered} Exam Preparation}}
\fancyfoot[L]{\tiny\color{black!55}\parbox[t]{0.78\textwidth}{AP\textsuperscript{\textregistered} is a registered trademark of the College Board, which was not involved in the production of, and does not endorse, this material. \textcopyright{} Integration Academy.}}
\fancyfoot[R]{\scriptsize\thepage}
\renewcommand{\headrulewidth}{1.2pt}
\renewcommand{\headrule}{\hbox to\headwidth{\color{gold}\leaders\hrule height \headrulewidth\hfill}}
\renewcommand{\footrulewidth}{0pt}
\setlength{\parindent}{0pt}
\setlength{\parskip}{4pt}

% Title block: \wstitle{code}{title}{unit line}
\newcommand{\wstitle}[3]{\begin{center}{\small\bfseries\color{gold}WORKSHEET #1}\\[3pt]{\LARGE\bfseries\color{navy}#2}\\[3pt]{\small\color{black!60}#3}\end{center}\vspace{2pt}}
% Section heading
\newcommand{\sech}[1]{\par\vspace{10pt}{\large\bfseries\color{navy}#1}\par\nobreak\vspace{2pt}{\color{rule}\hrule height0.8pt}\vspace{7pt}}
% A level tag printed beside a group of questions
\newcommand{\level}[1]{\par\vspace{6pt}{\footnotesize\bfseries\color{gold}\MakeUppercase{#1}}\par\nobreak\vspace{2pt}}

\newtcolorbox{ideabox}{colframe=ink,colback=inkbg,coltitle=white,colbacktitle=ink,fonttitle=\bfseries\small,title={Key ideas},arc=2pt,boxrule=0.8pt,left=8pt,right=8pt}
\newtcolorbox{example}[1]{colframe=rule,colback=white,coltitle=navy,colbacktitle=inkbg,fonttitle=\bfseries\small,title={#1},arc=2pt,boxrule=0.8pt,borderline west={3pt}{0pt}{ink},left=10pt,right=8pt}
\newtcolorbox{frbox}[1]{colframe=gold,colback=goldbg,coltitle=white,colbacktitle=gold,fonttitle=\bfseries\small,title={#1},arc=2pt,boxrule=0.8pt,breakable,left=8pt,right=8pt}
\newtcolorbox{solbox}{colframe=okgreen,colback=okbg,coltitle=white,colbacktitle=okgreen,fonttitle=\bfseries\small,title={Worked solutions},arc=2pt,boxrule=0.8pt,breakable,left=8pt,right=8pt}
\newtcolorbox{scorebox}[1]{colframe=navy,colback=white,coltitle=white,colbacktitle=navy,fonttitle=\bfseries\small,title={#1},arc=2pt,boxrule=0.8pt,breakable,left=8pt,right=8pt}

\newcommand{\soln}{\par\smallskip\textbf{\color{ink}Solution.}\ }
\newcommand{\result}[1]{\par\smallskip\textbf{#1}}
% blank working space of a given height
\newcommand{\work}[1]{\par\vspace{#1}}
% a point value in the margin of a scoring guide
\newcommand{\pts}[1]{\hfill{\footnotesize\bfseries\color{gold}[#1]}}

% Coordinate grid for graph questions: \begin{apgraph}{xmin}{xmax}{ymin}{ymax} ... \end{apgraph}
\newenvironment{apgraph}[4]{\begin{center}\begin{tikzpicture}\begin{axis}[
  axis lines=middle, xmin=#1, xmax=#2, ymin=#3, ymax=#4,
  grid=both, grid style={black!12}, xtick distance=1, ytick distance=1,
  ticklabel style={font=\scriptsize}, xlabel={$x$}, ylabel={$y$},
  width=10cm, height=7.2cm, samples=80, clip=false,
  every axis plot/.append style={line width=1.1pt, color=ink}]}{\end{axis}\end{tikzpicture}\end{center}}
\newcommand{\closeddot}[2]{\addplot[only marks, mark=*, mark size=2.2pt, color=ink] coordinates {(#1,#2)};}
\newcommand{\opendot}[2]{\addplot[only marks, mark=*, mark size=2.2pt, color=ink, mark options={fill=white}] coordinates {(#1,#2)};}

% Multiple-choice options: four across, or stacked when the options are long.
\newcommand{\opts}[4]{\par\vspace{3pt}\begin{tabular}{@{}p{0.235\linewidth}p{0.235\linewidth}p{0.235\linewidth}p{0.235\linewidth}@{}}(A)\ #1 & (B)\ #2 & (C)\ #3 & (D)\ #4\end{tabular}\par\vspace{9pt}}
\newcommand{\optsv}[4]{\par\vspace{3pt}\begin{tabular}{@{}l@{\ \ }p{0.9\linewidth}@{}}(A) & #1\\(B) & #2\\(C) & #3\\(D) & #4\end{tabular}\par\vspace{9pt}}

\fancyhead[R]{\footnotesize\color{black!60}Worksheet 0.4}
\begin{document}

\wstitle{0.4}{Your Graphing Calculator: The Four Exam Skills}{Unit 0: Readiness}

\begin{ideabox}
\begin{itemize}[leftmargin=*,itemsep=2pt,topsep=2pt]
  \item The exam expects four calculator skills: \textbf{graph} in a window, \textbf{solve} an equation, find a \textbf{derivative at a point}, and evaluate a \textbf{definite integral}.
  \item Work in \textbf{radian mode}. Give final answers correct to \textbf{three decimal places}.
  \item \textbf{Store} intermediate values at full precision; round once, at the end.
  \item Write the set-up in standard notation (an equation, $f'(a)$, or an integral with limits) \emph{before} the decimal. Calculator syntax earns nothing.
\end{itemize}
\end{ideabox}

\sech{Worked examples}
\begin{example}{Example 1 \textperiodcentered{} Solving an equation}
Find the positive solution of $x^4-x-3=0$.\soln Graph $y=x^4-x-3$; it crosses the positive $x$-axis once, between 1 and 2. Use the zero command.\result{$x^4-x-3=0$ when $x\approx1.453$.}
\end{example}

\begin{example}{Example 2 \textperiodcentered{} A derivative at a point}
Let $f(x)=x^2\ln x$. Find $f'(1.7)$.\soln Use the numerical derivative command at $x=1.7$.\result{$f'(1.7)\approx3.504$.}
\end{example}

\begin{example}{Example 3 \textperiodcentered{} A definite integral}
Evaluate $\displaystyle\int_1^2\sqrt{x^3+1}\,dx$.\soln Use the numerical integral command with limits 1 and 2.\result{$\displaystyle\int_1^2\sqrt{x^3+1}\,dx\approx2.130$.}
\end{example}

\newpage
\sech{Practice}
\emph{A graphing calculator is required for every question on this page. Work in radian mode and give decimal answers to three decimal places.}
\level{Core}
\begin{enumerate}[leftmargin=*,itemsep=2pt]
  \item Find the smaller of the two solutions of $e^{x}=4x$.\work{2.3cm}
  \item Find the value of $x$ at which the graphs of $y=x^3$ and $y=\cos x$ intersect.\work{2.3cm}
  \item Let $f(x)=\arctan\!\left(x^2\right)$. Find $f'(1.3)$.\work{2.3cm}
  \item Evaluate $\displaystyle\int_0^1e^{x^2}\,dx$.\work{2.3cm}
\end{enumerate}

\level{Exam level}
\begin{enumerate}[leftmargin=*,itemsep=2pt,start=5]
  \item Find the larger of the two solutions of $e^{x}=4x$.\work{2.8cm}
  \item The graphs of $y=4-x^2$ and $y=e^{x}$ intersect at $x=a$ and $x=b$, with $a<b$. Find $a$ and $b$, and find $\displaystyle\int_a^b\left(4-x^2-e^{x}\right)dx$.\work{2.8cm}
  \item Let $f(x)=\sin\!\left(e^{x}\right)$. Find $f'(0.5)$.\work{2.8cm}
  \item A student writes \textquotedblleft{}fnInt(X sin(X), X, 0, $\pi$) $=3.14$\textquotedblright{}. Rewrite this line so that it would earn full credit.\work{2.8cm}
\end{enumerate}

\begin{frbox}{Question 9 \textperiodcentered{} Free-response practice}
Let $f(x)=3\ln(x+1)$ and $g(x)=x^2$. The graphs of $f$ and $g$ intersect at $x=0$ and at $x=c$, where $c>0$.
\begin{enumerate}[label=(\alph*),leftmargin=*,itemsep=2pt]
  \item Find $c$. Write the equation you solved.\work{2.6cm}
  \item Find the area of the region enclosed by the two graphs. Write the integral before the decimal.\work{2.6cm}
  \item Find $f'(1)$ and $g'(1)$. Which graph is steeper at $x=1$?\work{2.6cm}
\end{enumerate}
\end{frbox}

\newpage
\begin{solbox}
\begin{enumerate}[leftmargin=*,itemsep=4pt]
  \item $e^{x}-4x=0$ when $x\approx0.357$.
  \item $x^3=\cos x$ when $x\approx0.865$.
  \item $f'(1.3)\approx0.674$.
  \item $\displaystyle\int_0^1e^{x^2}\,dx\approx1.463$.
  \item $x\approx2.153$.
  \item $a\approx-1.965$ and $b\approx1.058$ (stored at full precision). $\displaystyle\int_a^b\left(4-x^2-e^{x}\right)dx\approx6.428$.
  \item $f'(0.5)\approx-0.128$.
  \item $\displaystyle\int_0^{\pi}x\sin x\,dx\approx3.142$. Standard notation, and three decimal places.
  \item \begin{enumerate}[label=(\alph*),leftmargin=*,itemsep=3pt]
    \item $3\ln(c+1)=c^2$ gives $c\approx1.738$.
    \item $\displaystyle\int_0^{c}\left[3\ln(x+1)-x^2\right]dx\approx1.309$.
    \item $f'(1)=1.5$ and $g'(1)=2$, so the graph of $g$ is steeper at $x=1$.
  \end{enumerate}
\end{enumerate}
\end{solbox}

\end{document}
